% Objective review of the teacher's Question 2 proposal
\subsection{对教师方案的客观取舍}

教师方案提出的“缸内压力--Wiebe放热--$p\,\mathrm dV/\mathrm d\theta$气体转矩--摩擦--重叠角换气”链条是正确的物理主线，本文保留该结构。具体采纳三点改进：偏心距按附件尺寸取 $e=45-30=15\,\mathrm{mm}$；转子惯性力只作为瞬态轴承载荷，不计入完整循环的平均输出功率；用公开 AIE 225CS 台架曲线先标定 BMEP，再迁移到题设几何。

同时，本文对方案作出三点收敛。缸压和 Wiebe 参数需要实测缸压后辨识，不能由示意图唯一确定；重叠角不是换气质量流量，必须先通过端口面积、压差和流量系数积分得到无量纲吞吐量 $\xi$；在缺少喷油时刻和排放台架数据时，不能宣称存在唯一最优重叠角。因此本文对可辨识量给出中心值和区间，对不可辨识量给出基于全混合假设的灵敏度结果。这样既保留教师方案的机理闭环，又避免把设计先验误写成实验标定结论。

